\subsection*{Assignment 2 -- Data Cleaning}
L'Assignment 2 richiedeva di affrontare i problemi legati a valori
mancanti e incongruenze presenti nel dataset menzionate nella sezione
precedente.
Il processo è stato suddiviso in tre fasi principali:
\begin{itemize}
    \item \textbf{Data retrieval}: raccolta di informazioni mancanti da
    fonti esterne;
    \item \textbf{Integration}: unione dei dati esterni con quelli
    originali;
    \item \textbf{Fill}: completamento dei valori mancanti e
    correzione delle incongruenze.
\end{itemize}

Le funzioni relative a questi tre processi sono contenute all'interno
dei file omonimi nella cartella \texttt{assignment\_2}.

\subsubsection*{Data Retrieval}
Questa fase si basa sull'API \texttt{musicbrainzngs}.
Il processo divide il recupero di informazioni
relative agli \textbf{artisti} e alle \textbf{canzoni}, impiegando
comunque la stessa logica per le due entità.

A partire dai dati disponibili nel dataset, vengono costruite liste di
clausole formattate secondo la sintassi Lucene.
L'ordine di queste clausole è basato sull'affidabilità dell'attributo in
questione. Queste liste vengono poi utilizzate iterativamente, in ordine,
per generare query da inviare a Musicbrainz, finché non finiscono
le clausole, viene restituita una risposta vuota o una risposta con un
solo risultato.

Le query restituiscono risultati sempre più precisi; l'ordine delle clausole
nella lista è definito precisamente per questo motivo,
assicurando almeno un risultato per ogni artista o canzone.
Alla fine del processo, i dati ricavati sono stati salvati nei file della
cartella \texttt{dataset/mb}.


\subsubsection*{Integration}
Dopo il recupero dei dati da MusicBrainz, il passo successivo consiste
nell'integrazione delle informazioni ottenute con il dataset
originale.
Questo processo prevede il confronto tra gli elementi del dataset
originale e quelli recuperati, restituendo la lista dei risultati
più attinenti (uno per ogni artista o canzone),
con il proprio \texttt{matching\_score}
assegnato. Il calcolo del matching score è basato sulla somiglianza
degli attributi dell'oggetto originale e ricavato nel processo di
retrieval.
I valori di somiglianza vengono moltiplicati tra loro, in modo da
ottenere lo score finale compreso tra 0 e 1.
Questo approccio consente di integrare in maniera automatica e robusta
i dati esterni.


\subsubsection*{Fill}
L'ultima fase del processo riguarda il completamento dei valori mancanti,
utilizzando i dati del miglior risultato trovato precedentemente.  
Per ogni elemento del dataset originale, se la corrispondenza con i dati
recuperati è considerata sufficientemente affidabile
(matching score superiore a una soglia), vengono aggiornati gli
attributi mancanti.
Le informazioni geografiche vengono ulteriormente validate
tramite interrogazioni a OpenStreetMap, garantendo coerenza e
completezza dei dati. Per correggere l'incongruenza della presenza di
dati del luogo di nascita o della residenza, sono state scelte sempre
i luoghi di nascita, affidandosi ai dati ricavati da MusicBrainz.

\subsubsection*{Flusso Completo}

Il processo è successivamente stato implementato in un unico script.
Oltre alle operazioni precedenti, questo tratta anche:

\begin{enumerate}
    \item \textbf{Generazione di nuove identificazioni}: creazione di
    ID univoci tramite \texttt{uuid} per artisti, canzoni, album, date
    e località geografiche.
    
    \item \textbf{Costruzione delle tabelle}: dai file originali,
    lo script separa le entità delle date di pubblicazione
    delle canzoni, le informazioni sul luogo di nascita degli
    artisti, e le partecipazioni degli artisti alle canzoni. 
\end{enumerate}

Questo flusso garantisce un dataset finale completo, coerente e pronto
per le analisi successive, con tutte le informazioni chiave normalizzate
e integrate.
